%\documentclass[10pt,respuestas,a4]{examenEPN}
\documentclass[10pt,a4]{examenEPN}

\usepackage{comandosEPN}
\usepackage{esint}
\usepackage{tikz}
\usepackage{systeme}
\usetikzlibrary{automata, positioning}
\usepackage{multicol}

\usepackage{pstricks,pst-plot}

\input{Comandos.tex}


\materia{Álgebra Lineal}
\examen{Repaso previo al examen bimestral}
\autor{Departamento de Formación Básica}
\fecha{Semestre 2024-B}
\logouno[1.6cm]{Logo_EPN.eps}
\logodos[2.6cm]{Logo_DFB.eps}

\begin{document}

\encabezado
\begin{preguntas}
%%%%%%%%%%%%%%%%%%%%%%%%%%%%%%%%%%%%%%%%%%%%%%%%
\item Dada la aplicación lineal:
\[
    \funcion{T}{\R^4}{\R^{2\times 2}}{(a,b,c,d)}{\left( \begin{matrix}
        3a-b & c+d \\
        2c-a & a \\
     \end{matrix}\right)}
\]
Considere la base de $\mathbb{R}^4$,

$ B=\left\{(3,2,2,-2), (0,0,-3,2),(0,2,-3,1),(0,3,-1,-2)   \right\}$


Encuentre la matriz asociada $[T]_C^B$. Donde $C$ es la base can\'onica. ¿Esta matriz es invertible? Si lo es, indique qu\'e le dice sobre la aplicación $T$.
\begin{respuesta}
La matriz $[T]_C^B$ es invertible. Eso se lo puede ver desde inicio al comprobar que $T$ es biyectiva
\end{respuesta}

%%%%%%%%%%%%%%%%%%%%%%%%%%%%%%%%%%%%%%%%%%%%%%%%
\item Dada la aplicación lineal:
\[
    \funcion{H}{\R^{2\times 2}}{\R^4}{\left( \begin{matrix}
        a & b \\
        c & d \\
     \end{matrix}\right)}{(a,c,b,d)}
\]

Considere la aplicaci\'on $T$ y las bases de los espacios del ejercicio anterior, y encuentre la matriz asociada de $H$, $(A=[H]_B^C)$.


%%%%%%%%%%%%%%%%%%%%%%%%%%%%%%%%%%%%%%%%%%%%%%%%
\item Considerando nuevamente la aplicaci\'on $H$ del ejercicio 2, obtenga una matriz diagonal que sea semejante a $[H]_{C_2}^{C_1}$. Donde $C_1$ y $C_2$ son las bases can\'onicas respectivas.
\begin{respuesta}
Diagonalizar la matriz respectiva
\end{respuesta}

%%%%%%%%%%%%%%%%%%%%%%%%%%%%%%%%%%%%%%%%%%%%%%%%
\item Considere la aplicaci\'on $T\in\mathscr{L}(\mathbb{R}^3,\mathbb{R}^3)$ y $u,v \in \mathbb{R}^3$ tal que $Tu=2u \text{ y } T(v)=3v$. Si se conoce que $T(0,1,0)=(0,\pi,0)$, ¿el conjunto 
$B=\left\{u,v,w\right\}$, donde $w=(0,e,0)$, es una base de $\mathbb{R}^3$?
\begin{respuesta}
Diagonalizar la matriz respectiva
\end{respuesta}
%%%%%%%%%%%%%%%%%%%%%%%%%%%%%%%%%%%%%%%%%%%%%%%%%%%%%%%%%%%%%%%%%%%%%%%%%%%%%%%
\item Diagonalice la matriz
\begin{center}
 $A = \left( \begin{matrix}
        8 & 0 & -10\\
        2 & 1 & -2\\
        5 & 0 &  -7\\
     \end{matrix}\right)$.   
\end{center}
\begin{respuesta}
Los valores propios de la matriz son $\lambda_1=1, \lambda_2=3 \text{y} \lambda_3=-2 $

\end{respuesta}
%%%%%%%%%%%%%%%%%%%%%%%%%%%%%%%%%%%%%%%%%%%%%%%%%%%%%%%%%%%%%%%%%%
%%%%%%%%%%%%%3
\item Considere la siguiente matriz 
\begin{center}
 $A = \left( \begin{matrix}
        1 & 1 & 3\\
        3 & 1 & -5\\
        -1 & 1 &  5\\
     \end{matrix}\right)$  
\end{center}
Uno de sus valores propios es $3$. Indique si esta matriz es diagonalizable justificando su respuesta.
\begin{respuesta}
$A$ no es diagonalizable. Para la matriz $A$, La $\dim(E_2)$ es estrictamente menor que la multiplicidad algebraica del valor propio $2$. 
\end{respuesta}
%%%%%%%%%%%%%%%%%%%%%%%%%%%%%%%%%%%%%%%%%%%%%%%%%%%%%%%%%%%%%%%%%%%%%%%%%%%%%%%4
\item Diagonalice la matriz
\begin{center}
 $A = \left( \begin{matrix}
        1 & 0 & -1\\
        0 & 1 & 2\\
        -1 & 2 &  5\\
     \end{matrix}\right)$.   
\end{center}
\begin{respuesta}
Los valores propios de la matriz son $\lambda_1=0 \text{y} \lambda_2=1 \text{ y} \lambda_2=6$.

\end{respuesta}
%%%%%%%%%%%%%%%%%%%%%%%%%%%%%%%%%%%%%%%%%%%%%%%%%%%%%%%%%%%%%%%%%%%%%%%%%%%%%%%4
\item Diagonalice la matriz
\begin{center}
 $A = \left( \begin{matrix}
        -3 & 1 & 2\\
        1 & -3 & 2\\
        2 & 2 &  0\\
     \end{matrix}\right)$.   
\end{center}
\begin{respuesta}
Los valores propios de la matriz son $\lambda_1=1 \text{y} \lambda_2=-4$. Donde el \'ultimo valor propio tiene multiplicidad $2$.

\end{respuesta}
%%%%%%%%%%%%%%%%%%%%%%%%%%%%%%%%%%%%%%%%%%%%%%%%%%%%%%%%%%%%%%%%%%
%%%%%%%%%%%%%3

\item Dada la aplicación lineal:
\[
    \funcion{T}{\R^3}{\R^3}{(a,b,c)}{(-4a+3b-6c,b,3a-2b+5c)}
\]
\begin{enumerate}
    \item Determine los valores propios de $T$.
    \item 
    determine los vectores propios de la matriz $A$ y de la aplicación $T$.
    \item 
    Determine una matriz $D$ diagonal, semejante a la matriz $A$.
    \item 
    Determine una matriz $Q$, tal que $QD=AQ$
\end{enumerate}
\begin{respuesta}
Sugerencia: encuentre la matriz asociada en la base can\'onica y diagonal\'icela. Luego, encuentre la base en la que la matriz asociada de $T$ es la diagonal encontrada anteriormente. 
\end{respuesta}
\item 
Diagonalice ortogonalmente la matriz
\begin{center}
 $A = \left( \begin{matrix}
        0 & 1 & 1\\
        1 & 0 & 1\\
        1 & 1 &  0\\
     \end{matrix}\right)$.   
\end{center}
\end{preguntas}
\final
\end{document}
